% English translation of checked_pages/esq_119_018.tex; source: cote 119, site 18, PDF page 19, typescript page 16.
% AI-translated by GPT-6 Astra (OpenAI, ChatGPT Pro mode).
\scanpage{119/18}
\mspage{16}

that the central theme of anabelian algebraic geometry appeared:
reconstructing certain varieties $X$ called ``anabelian''
over an absolute field $K$ from their mixed fundamental group,
an extension of $\mathrm{Gal}(\bar{K}/K)$ by $\pi_1(X_{\bar{K}})$; it was then that the
``fundamental conjecture of anabelian algebraic geometry'' emerged,
close to the conjectures of Mordell and Tate that Faltings
has just proved($^3$). It was also there that a first reflection on
Teichmüller groups and the first intuitions about the multiple
structure of the ``Teichmüller tower'' began - the open moduli
multiplicities $M_{g,\nu}$ appearing, moreover, as the first important
examples, in dimension $> 1$, of varieties (or rather,
multiplicities) that seem well deserving of the name ``anabelian''.
Toward the end of this period of reflection, it appeared to me
as a foundational reflection on a theory still in
limbo, for which the name ``\uline{Galois-Teichmüller theory}''
seemed more appropriate than ``Galois Theory'', the name I had initially
given to my notes.

This is not the place to give a more detailed account
of this whole collection of questions, intuitions, ideas - including
tangible results, certainly. The most important seems to me to be the one
mentioned in passing in par. 2, namely the faithfulness of the outer
action of $\mathrm{I}\!\Gamma = \mathrm{Gal}(\bar{\mathbb{Q}}/\mathbb{Q})$ (and of its open subgroups) on $\hat{\pi}_{0,3}$,
and more generally (if I remember correctly) on the fundamental
group of any ``anabelian'' algebraic curve (i.e. whose genus
$g$ and ``number of holes'' $\nu$ satisfy the inequality $2g + \nu \geq 3$,
i.e. such that $\chi(X) < 0$) defined over a finite extension of $\mathbb{Q}$. This
result can be regarded as essentially equivalent to
Bielyi's theorem - it is the first concrete manifestation, through
a precise mathematical statement, of the ``message'' discussed
above.

I would like to finish this rapid overview with a few words
of commentary on the truly unimaginable richness of a typical
anabelian group such as the group $S\ell(2,\mathbb{Z})$ - probably the most
remarkable infinite discrete group one has encountered, which appears

\margin{18}

